% SOLUCIÓN Ejercicio 4 — documento corregido
\documentclass[12pt,a4paper]{article}
\usepackage[utf8]{inputenc}
\usepackage[T1]{fontenc}
\usepackage[spanish]{babel}
\usepackage{amsmath,amssymb}
\usepackage{booktabs}
\usepackage{siunitx}
\sisetup{locale=FR, output-decimal-marker={,}}
\usepackage{circuitikz}
\usepackage[hidelinks]{hyperref}

\title{Ley de Ohm}
\author{Estudiante}
\date{\today}
\begin{document}
\maketitle
\tableofcontents

\section{Introducción}
La ley de Ohm establece que
\begin{equation}
  V = I\cdot R
  \label{eq:ohm}
\end{equation}
Para dos resistencias en serie, por la ecuación~\eqref{eq:ohm} y la LVK,
\begin{equation}
  R_{eq} = R_1 + R_2
\end{equation}

\section{Resultados}
Los resultados se ven en la Tabla~\ref{tab:ohm} y en la Figura~\ref{fig:ohm}.

\begin{table}[htbp]
  \centering
  \caption{Mediciones de voltaje y corriente}
  \label{tab:ohm}
  \begin{tabular}{S[table-format=1.1] S[table-format=2.1]}
    \toprule
    {$V$ (\si{\volt})} & {$I$ (\si{\milli\ampere})} \\
    \midrule
    1,0 & 9,8 \\
    5,0 & 49,2 \\
    \bottomrule
  \end{tabular}
\end{table}

\begin{figure}[htbp]
  \centering
  \fbox{\parbox{0.6\textwidth}{\centering\vspace{1.5cm}Figura aquí\vspace{1.5cm}}}
  \caption{Montaje experimental}
  \label{fig:ohm}
\end{figure}
\end{document}

% Cambios (ver soluciones/README.md): babel español y fontenc; «Introducción» con tilde;
% \\ vacíos eliminados; $$ -> equation; «V = I x R» -> $V=I\cdot R$ numerada;
% referencias manuales -> \ref con label; tabla sin líneas verticales (booktabs),
% con caption/label; decimales con coma y siunitx; figura con caption/label.
